\documentclass[12pt, spanish]{article}
\usepackage[spanish]{babel}
\selectlanguage{spanish}
%\usepackage{natbib}
\usepackage{url}
\usepackage[utf8x]{inputenc}
\usepackage{graphicx}
\graphicspath{{images/}}
\usepackage{parskip}
\usepackage{fancyhdr}
\usepackage{vmargin}
\usepackage{multirow}
\usepackage{float}
\usepackage{chngpage}
\usepackage{enumitem}
\usepackage{forloop}


\usepackage{amsfonts}

\usepackage{subcaption}

\usepackage{hyperref}
\usepackage[
    type={CC},
    modifier={by-nc-sa},
    version={4.0},
]{doclicense}

\hypersetup{
    colorlinks=true,
    linkcolor=blue,
    filecolor=magenta,
    urlcolor=cyan,
}

% para codigo
\usepackage{listings}
\usepackage{xcolor}


\usepackage[default]{sourcesanspro}

\setmarginsrb{2 cm}{1 cm}{2 cm}{2 cm}{1 cm}{1.5 cm}{1 cm}{1.5 cm}

\title{Práctica 1:\\
Aritmética modular.\hspace{0.05cm} }
\author{Antonio David Villegas Yeguas}
\date{\today}

\renewcommand*\contentsname{hola}

\makeatletter
\let\thetitle\@title
\let\theauthor\@author
\let\thedate\@date
\makeatother

\pagestyle{fancy}
\fancyhf{}
\rhead{\theauthor}
\lhead{\thetitle}
\cfoot{\thepage}


\usepackage{caption}

\lstset{
language=C++,
basicstyle=\small\ttfamily,
numbers=left,
numbersep=5pt,
xleftmargin=20pt,
frame=tb,
framexleftmargin=20pt
}

\renewcommand*\thelstnumber{\arabic{lstnumber}:}

\DeclareCaptionFormat{mylst}{\hrule#1#2#3}
\captionsetup[lstlisting]{format=mylst,labelfont=bf,singlelinecheck=off,labelsep=space}

\begin{document}

%%%%%%%%%%%%%%%%%%%%%%%%%%%%%%%%%%%%%%%%%%%%%%%%%%%%%%%%%%%%%%%%%%%%%%%%%%%%%%%%%%%%%%%%%

\begin{titlepage}
    \centering
    \vspace*{0.3 cm}
    \includegraphics[scale = 0.50]{ugr.png}\\[0.7 cm]
    %\textsc{\LARGE Universidad de Granada}\\[2.0 cm]
    \textsc{\large 4º CSI 2020/21}\\[0.5 cm]
    \textsc{\large Grado en Ingeniería Informática}\\[0.5 cm]
    \rule{\linewidth}{0.2 mm} \\[0.2 cm]
    { \huge \bfseries \thetitle}\\
    \rule{\linewidth}{0.2 mm} \\[1 cm]

    \begin{minipage}{0.4\textwidth}
        \begin{flushleft} \large
            \emph{Autor:}\\
            \theauthor\\
			 \emph{DNI:}\\
            77021623-M
            \end{flushleft}
            \end{minipage}~
            \begin{minipage}{0.4\textwidth}
            \begin{flushright} \large
            \emph{Asignatura: \\
            Criptografía}   \\
            \emph{Correo:}\\
            advy99@correo.ugr.es
        \end{flushright}
    \end{minipage}\\[0.5cm]

    {\large \thedate}\\[0.5cm]
    {\url{https://github.com/advy99/CRIP/}}
    {\doclicenseThis}

    \vfill

\end{titlepage}

%%%%%%%%%%%%%%%%%%%%%%%%%%%%%%%%%%%%%%%%%%%%%%%%%%%%%%%%%%%%%%%%%%%%%%%%%%%%%%%%%%%%%%%%%

\tableofcontents
\pagebreak

%%%%%%%%%%%%%%%%%%%%%%%%%%%%%%%%%%%%%%%%%%%%%%%%%%%%%%%%%%%%%%%%%%%%%%%%%%%%%%%%%%%%%%%%%


\section*{Introducción}

En esta práctica repasaremos conceptos de las asignaturas Álgebra Lineal y Métodos Discretos y Lógica y Métodos Discretos de cara a implementar algoritmos que nos permitan trabajar en anillos y cuerpos de $\mathbb{Z}$.

\section{Algoritmo extendido de Euclides}

En este primer ejercicio se nos pide implementar el algoritmo extendido de Euclides. Este algoritmo nos permitirá obtener el máximo común divisor de dos números que recibirá como entrada.

Este algoritmo se basa en ir realizando divisiones sucesivas hasta obtener un cero, y el último resto distinto de cero será el máximo común divisor. Además, también obtendremos un $s$ y un $t$, tal que si buscamos el máximo común divisor de $a$ y $b$, $mcd(a, b) = a \cdot s + b \cdot t$.

La implementación es la siguiente:

\begin{lstlisting}[caption={Algoritmo extendido de Euclides}]
INPUT: a, b
OUTPUT: r <- mcd(a,b), s, t

Si a == 0 o b == 0
	Devolver vacio

r0 <- a
r1 <- b
s0 <- 1
s1 <- 0
t0 <- 0
t1 <- 1

Mientras r1 != 0:
	q <- r0 / r1
	r0, r1 <- r1, r0 - r1 * q
	s0, s1 <- s1, s0 - s1 * q
	t0, t1 <- t1, t0 - t1 * q

Devolver r0, s0, t0
\end{lstlisting}

Como vemos, el algoritmo se basa en ir dividiendo $a$ por $b$, y actualizar sus valores, $a$ pasa a valer $b$, y $b$ el residuo de la división. De esta forma, cuando $b$ llegue a cero, $a$ valdrá el residuo anterior, que será el máximo común divisor.

En $s$ y $t$ iremos almacenando los valores de la identidad de Bezout de cara a tener los coeficientes con los que se puede expresar el máximo común divisor.


\section{Inverso modular de dos primos relativos}

Para este ejercicio vamos a calcular el inverso de un número $a$ módulo $b$, siendo $a$ y $b$ dos primos relativos.

Si $a$ y $b$ son primos resativos, su máximo común divisor vale uno, por lo que sabemos por la identidad de Bezout que:

$$ mcd(a, b) = 1 = u \cdot a + v \cdot b $$

Y también sabemos que estamos trabajando en los restos de dividir por $b$, es decir $\mathbb{Z}_b$, luego:

$$1 = u \cdot a + v \cdot b $$
$$1 = u \cdot a $$
$$a^{-1} = u$$

Ya que en $\mathbb{Z}_b$ el valor de $b$ es cero. De esta forma, con el algoritmo extendido de Euclides y sabiendo que $a$ y $b$ son primos relativos podemos obtener el inverso de $a$ módulo $b$.

En este caso utilizamos el algoritmo del ejercicio anterior, usando el segundo valor devuelto, el de $s$. También podríamos volverlo a escribir sin calcular $t$, ya que no es necesario para el cálculo.


\section{Potencia modular}

Ahora se nos pide implementar un algoritmo de forma que calculemos $a^b \mod n$ para cualquier $a$, $b$ y $n$ enteros positivos.

Este ejercicio parece simple a primera vista, pero el problema es que si se tratan de valores excesivamente grandes la operación normal de exponenciación no es suficiente, ya que desbordaría los registros y memoria de la máquina.

Para evitar esto, sabiendo que vamos a trabajar en $\mathbb{Z}_{n}$, la idea es aprovechar la representación binaria del exponente para saber cuando tenemos que multiplicar el resultado por la base, y al utilizar representación binaria, cada vez que desplacemos un bit del exponente, elevar al cuadrado la base, de forma que la operación sea equivalente.


\begin{lstlisting}[caption={Algoritmo de exponenciación modular}]
INPUT: base, exponente, modulo
OUTPUT: resultado <- (base^exponente) mod modulo

resultado <- 1

Mientras exponente > 0:
	Si el bit menos significativo del exponente == 1:
		resultado <- (resultado * base) mod modulo

	exponente <- exponente desplazando su valor un bit a la derecha
	base <- (base * base) mod modulo

Devolver resultado
\end{lstlisting}


\section{Test de Miller-Rabin}

\section{Logaritmo discreto}

\section{Raices cuadradas modulares}

\subsection{Raiz cuadrada modular}

\subsection{Obtener todas las raices}


\section{Métodos de factorización}

\subsection{Factorización utilizando el método de Fermat}

\subsection{Factorización utilizando el algoritmo de $\rho$ de Pollard}



\section{Comparación de tiempos}


% \begin{thebibliography}{9}
%
%
% \end{thebibliography}

\end{document}
